\documentclass[instructions.tex]{subfiles}
\begin{document}
 
\chapter{Du jugement.}
 
Le jour du jugement est appelé dans les livres saints, le jour du Seigneur, parce qu'alors Dieu se fera connoître, en manifestant avec plus de gloire sa grandeur et sa puissance. Ce jour sera annoncé par d'horribles tremblemens de terre, par l'éclipse de tous les astres, par l'embrasement de l'univers, par la résurrection des morts, et par des prodiges si effrayans, que les hommes en sécheront de crainte\footnote{Arescentibus hominibus præ timore. \textit{Luc.} 21. 26.}.
 
Quatre circonstances seront insupportables aux pécheurs dans ce jugement : la présence de Jésus-Christ, le compte qu'il faudra rendre, leur confusion devant l'univers, la sentence de réprobation.
 
\primo La vue et la présence de Jésus-Christ au jugement sera plus insupportable au pécheur que les feux de l'enfer. Si le Sauveur, lorsqu'il vivoit sur la terre, n'avoit pas caché l'éclat de sa majesté, aucun mortel n'auroit pu soutenir la présence de cet Homme-Dieu. Les rois Mages, en le voyant sous la forme d'un enfant, furent remplis d'un respect si profond, qu'ils se prosternèrent pour l'adorer. Les apôtres ne virent sur le mont Thabor que quelques rayons de sa majesté; ils en furent si frappés, qu'ils tombèrent par terre. Le prophète Daniel, à la vue d'un ange, fut saisi d'une si grande frayeur, qu'il tomba la face contre terre. O Dieu! de quelle frayeur seront donc saisis les pécheurs, lorsqu'ils verront celui devant qui les anges mêmes tremblent de respect; lorsqu'ils le verront, non pas comme les Mages, sous la figure d'un enfant, mais sur le trône de sa justice, portant dans ses regards les marques d'un implacable courroux!
 
Mais, en voyant leur juge, hélas! que verront-ils? Ils verront Jésus-Christ, ils verront celui qu'ils ont crucifié\footnote{Videbunt in quem transfixerunt. \textit{Joan.} 19. 37.}; ils le verront précédé de sa croix. Alors, dit l'Évangile, toutes les nations pousseront des cris lamentables, à la vue d'un Dieu qui vouloit les sauver et qui vient les juger. Les pécheurs chercheront les cavernes pour se cacher, conjurant les rochers de tomber sur eux, pour les dérober à la vue de leur juge\footnote{Montes, cadite super nos. \textit{Luc.} 33. 30.}. Rassurez-vous, gens de bien, vous verrez votre juge avec consolation; tandis que les pécheurs, aujourd'hui si fiers, frémiront en sa présence.
 
II. La seconde circonstance qui accablera les pécheurs, c'est le compte qu'ils rendront de leur vie. Rien n'échappera à la connoissance du juge souverain; toutes les actions, les paroles, les pensées y seront examinées dans toute leur étendue de bonté ou de malice, et y recevront leur récompense ou leur châtiment propre.
 
Quelle sera votre surprise à la vue de tant de pensées et de désirs sensuels, de tant d'actions impures, de tant de paroles dissolues dans leurs expressions et dans leurs sens, de tant de vanité et d'orgueil, et de tant d'injustices et d'attachement aux biens de la terre, de tant d'emportemens, de colères et de rancunes, de tant de grâces négligées et de temps perdu! Vous serez effrayé du compte d'un seul jour; vous en serez même si accablé, que vous serez insupportable à vous-même. Eh mon Dieu, que sera-ce du compte de tant d'années, et du compte de toute la vie? Je suis assuré, disoit un saint homme, que le reproche que Dieu fera aux pécheurs réprouvés, d'une heure mal employée, sera le moment le plus insupportable de leur éternité. Quelle sera donc ma confusion, après avoir perdu des millions d'heures!
 
Job, cet homme juste, trembloit en pensant qu'il seroit jugé par un Dieu : Que deviendrai-je quand Dieu se lèvera pour me juger? Que lui dirai-je quand il m'aura examiné et convaincu? C'est ce qui le faisoit craindre dans toutes ses actions, parce qu'il les faisoit devant un juge qui sera alors inexorable\footnote{Verebar omnia opera mea, sciens quòd non parceres delinquenti. \textit{Job.} 9. 28.}.
 
Pourquoi ce saint homme craignait-il? Sa conscience ne lui rendoit-elle pas témoignage de son innocence? Ah! c'est qu'il savoit que Dieu juge les justes, et que la sainteté la plus épurée disparaît devant lui, comme la lumière d'un flambeau disparaît en présence du soleil. Hélas! disoit-il, en me croyant innocent, je serai trouvé criminel\footnote{Si innocentem me ostendero, pravum me comprobabit. 9. 20.}.
 
Malheur même à la vie la plus sainte, si vous l'examinez sans miséricorde, ô mon Dieu! disoit saint Augustin\footnote{Væ etiam laudabili vitæ hominum, si, remotâ misericordiâ, discutias eam. S. \textit{Aug.}.}. Un saint religieux disoit, en versant des larmes, à ceux qui le rassuraient à la mort : Oh! qu'il est terrible d'être obligé de rendre compte à Dieu! Hélas! dit saint Pierre, si le juste sera à peine sauvé, que deviendront le pécheur et l'impie\footnote{Si justus vix salvabitur, impius et peccator ubi parebunt? 1. \textit{Petr.} 4. 18.}? Un moyen assuré pour paroître avec confiance au jugement de Dieu, c'est de le prévenir en se jugeant soi-même; c'est de pleurer et de punir les péchés passés, par la pénitence, et de s'abstenir d'en faire de nouveaux.
 
\section{Résolutions}
 
\primo Je m'efforcerai d'apaiser mon redoutable juge, en lui demandant pardon du passé, tous les jours de ma vie.
 
\secundo Et je me concilierai un jugement favorable, en fuyant désormais le péché, et en remplissant tous mes devoirs avec ferveur et exactitude, dans la vue de plaire à Dieu, et de satisfaire à sa justice par les mérites de Jésus-Christ. 

\prayer{O mon Dieu! oubliez les nombreuses transgressions dont je me suis rendu coupable jusqu'ici, et faites-moi la grâce de ne plus rien faire qui mérite châtiment.}
 
\end{document}
